%% kapa: exam=1 round=19 year=2008 period=1 form=A subject=appraisal_law subject_ko=부동산관계법규 seq=18 number=58
%% source: Q-Net 감정평가사 기출문제 (cst00309) artlSeq=1001962 file=extracted/1008/1001962/2008년 제19회 감정평가사 1차시험 기출문제/1교시.hwp
%% conversion: kordoc 4.24.51 -> tools/gen1.py (auto); review=proofread-sonnet; flags=none
\begin{question}{58}
부동산 가격공시 및 감정평가에 관한 법령에 의할 때 ‘분할․합병등이 발생한 토지’에 해당하지 아니하는 것은?
\begin{choices}
  \item 지적법상 분할 또는 합병된 토지
  \item 하천법상 하천구역으로 결정된 토지
  \item 공유수면매립 등으로 지적법상 신규등록된 토지
  \item 토지의 형질변경 또는 용도변경으로 지적법상 지목변경이 된 토지
  \item 국·공유지가 매각 등의 사유로 사유지로 된 토지로서 개별공시지가가 없는 토지
\end{choices}
\end{question}
\answer{2}
